% ─── BASE TEÓRICA ─────────────────────────────────────────
\section{Base Teórica}

Esta seção apresenta a fundamentação teórica que sustenta o
trabalho, organizada de modo a evidenciar a lógica de
construção da pesquisa: o problema organizacional é
compreendido a partir da Engenharia de Processos, o
redesenho do fluxo é conduzido segundo a disciplina do
\textit{Business Process Management} (BPM), a transformação
desse redesenho em solução computacional é orientada pela
Engenharia de Software e pela Engenharia de Requisitos, a
implementação é viabilizada pela abordagem \textit{low-code},
e o resultado operacional é analisado à luz da automação de
processos. Dessa forma, a tecnologia não é apresentada como
ponto de partida do trabalho, mas como consequência de um
processo de análise e redesenho previamente conduzido.

\subsection{Engenharia de Processos e Transformação
Organizacional}

A Engenharia de Processos constitui a etapa inicial da
transformação digital, na medida em que estabelece que a
implantação de qualquer solução tecnológica deve decorrer de
uma necessidade operacional previamente identificada, e não
o contrário. Sob essa perspectiva, compreender o processo
existente --- seus participantes, atividades manuais, pontos
de falha e oportunidades de melhoria --- é condição prévia
para que a automação produza ganhos reais de padronização,
rastreabilidade e governança \cite{SantosCaldeiraPacheco2026}.

Um dos instrumentos centrais dessa abordagem é a
representação do processo em dois estados complementares.
O processo atual, ou \textbf{AS-IS}, descreve o fluxo antes
da intervenção, evidenciando comunicação descentralizada,
ausência de padronização, dependência de controles manuais e
baixa rastreabilidade das informações. O processo futuro, ou
\textbf{TO-BE}, representa o fluxo redesenhado após a
implantação da solução, contemplando abertura estruturada,
registro centralizado, definição clara de responsáveis,
automação de notificações e acompanhamento do ciclo de vida
de cada ocorrência \cite{SantosCaldeiraPacheco2026}. A
transição entre esses dois estados não decorre da simples
transposição do fluxo manual para um ambiente digital: quando
a digitalização ocorre sem revisão estrutural do processo, ela
tende a reproduzir, em formato eletrônico, as mesmas
ineficiências do modelo anterior
\cite{SantosCaldeiraPacheco2026}. O foco desta pesquisa,
portanto, não é demonstrar que um novo sistema foi criado,
mas evidenciar que o processo de gestão de chamados internos
foi analisado e redesenhado antes de sua automação.

\subsection{Gestão de Processos de Negócio (BPM)}

O \textit{Business Process Management} (BPM) fornece o
arcabouço metodológico que estrutura o ciclo de
identificação, modelagem, análise, redesenho, automação e
monitoramento de um processo organizacional
\cite{Ramos2024}. Trata-se de uma disciplina que orienta a
organização a superar a lógica funcional tradicional,
substituindo-a por uma visão orientada a processos, capaz de
atravessar diferentes áreas e gerar valor de ponta a ponta
\cite{SantosCaldeiraPacheco2026}.

No contexto de chamados internos entre áreas, o BPM orienta a
definição dos estados possíveis do chamado, as
responsabilidades de cada ator no fluxo e os critérios de
encerramento. \cite{Lopes2024} demonstraram que a
digitalização de processos com BPM e ferramentas do
ecossistema Microsoft resultou em repositório único de dados
com rastreabilidade completa, eliminando a comunicação
informal entre setores. Em sentido semelhante,
\cite{Ramos2024} evidenciam que a disponibilização de dados
gerenciais estruturados constitui o resultado de maior valor
percebido após iniciativas de BPM em instituições de médio e
grande porte, ao passo que
\cite{SantosCaldeiraPacheco2026} destacam que a efetividade
da transformação digital depende da revisão estruturada dos
processos e do fortalecimento das capacidades
organizacionais, e não apenas da incorporação de novas
ferramentas.

O projeto desenvolvido neste trabalho é relacionado às
etapas do ciclo de BPM conforme apresentado no
Quadro~\ref{quadro:bpm-projeto}.

\begin{table}[H]
\caption{Relação entre as etapas do BPM e o projeto desenvolvido}
\label{quadro:bpm-projeto}
\centering
\begin{tabular}{p{2.6cm} p{9.2cm}}
\hline
\textbf{Etapa do BPM} & \textbf{Aplicação no projeto} \\
\hline
Identificação & Análise do processo de comunicação de
chamados internos entre áreas. \\
\hline
Modelagem & Construção da representação do fluxo atual
(AS-IS) e do fluxo futuro (TO-BE). \\
\hline
Análise & Identificação dos gargalos operacionais a partir
do levantamento de requisitos junto aos participantes. \\
\hline
Redesenho & Definição do novo fluxo de atendimento, com
estados, responsáveis e regras de negócio. \\
\hline
Automação & Desenvolvimento dos componentes no ecossistema
Microsoft 365 (Power Apps, Microsoft List e Power
Automate). \\
\hline
Monitoramento & Construção do painel gerencial no Power BI
como base para indicadores de desempenho. \\
\hline
\end{tabular}
\raggedright
\par
{\fontsize{10}{12}\selectfont Fonte: Elaborado pelo autor,
adaptado de \cite{Ramos2024} e \cite{SantosCaldeiraPacheco2026}}
\end{table}

A utilização do BPM neste trabalho reforça que a tecnologia
foi aplicada como consequência da melhoria do processo, e
não como ponto inicial da solução.

\subsection{Engenharia de Software em Ambiente Low-Code}

A Engenharia de Software é a disciplina responsável por
transformar necessidades organizacionais em uma solução
tecnológica estruturada, por meio da aplicação sistemática de
métodos de análise, projeto, implementação e validação
\cite{Alves2022}. A adoção de uma plataforma \textit{low-code}
não elimina essa disciplina: ela modifica a forma de
implementação, substituindo a codificação manual extensiva
por configuração visual e componentes pré-construídos, mas
mantém a necessidade de análise de requisitos, definição de
arquitetura, estruturação de dados e validação do
comportamento esperado \cite{Mendes2024}.

Este trabalho relaciona explicitamente as etapas tradicionais
da Engenharia de Software às atividades desenvolvidas em
ambiente \textit{low-code}: o levantamento da necessidade
corresponde à pesquisa de campo conduzida junto aos três
perfis profissionais envolvidos no processo; a especificação
de requisitos corresponde à consolidação dos requisitos
funcionais e não funcionais; a definição de arquitetura
corresponde à escolha dos componentes do ecossistema
Microsoft 365 e à forma como eles se integram; a
implementação corresponde à construção do aplicativo, da
lista e dos fluxos de automação; e os testes e a validação
correspondem à verificação do funcionamento integrado e à
avaliação da solução junto aos participantes. Essa
correspondência evidencia que, mesmo em ambiente
\textit{low-code}, o projeto preserva as características de
um trabalho de engenharia, com requisitos funcionais e não
funcionais bem definidos.

\subsection{Engenharia de Requisitos}

A Engenharia de Requisitos é a área da Engenharia de Software
responsável por identificar, documentar, validar e gerenciar
as necessidades dos \textit{stakeholders} ao longo do ciclo
de desenvolvimento de um sistema. A qualidade dos requisitos
influencia diretamente a qualidade da solução produzida, pois
falhas introduzidas nessa etapa tendem a gerar retrabalho e
aumento de custos \cite{Alves2022}. Segundo
\cite{CarvalhoJunior2022}, a elicitação constitui uma das
fases mais críticas da Engenharia de Requisitos, pois a
qualidade de um software depende da correta compreensão das
necessidades dos usuários e da capacidade de alinhar os
requisitos gerados às metas do negócio --- alinhamento que,
segundo os mesmos autores, é favorecido quando os requisitos
são derivados a partir da modelagem prévia do processo de
negócio, aproximando a Engenharia de Requisitos da Engenharia
de Processos discutida anteriormente.

Neste trabalho, cada necessidade identificada junto aos
participantes da pesquisa foi conduzida por um percurso
específico até se tornar parte da solução, conforme
ilustrado no exemplo do Quadro~\ref{quadro:necessidade-requisito}.

\begin{table}[H]
\caption{Do problema organizacional ao requisito
implementado}
\label{quadro:necessidade-requisito}
\centering
\begin{tabular}{p{3.4cm} p{3.9cm} p{4.5cm}}
\hline
\textbf{Necessidade identificada} & \textbf{Requisito
funcional} & \textbf{Implementação} \\
\hline
Dificuldade para acompanhar quem está responsável pelo
chamado. & O sistema deve permitir identificar o responsável
pela tratativa. & Campo de responsável pela tratativa
integrado ao fluxo de atendimento na Microsoft List. \\
\hline
\end{tabular}
\raggedright
\par
{\fontsize{10}{12}\selectfont Fonte: Elaborado pelo autor}
\end{table}

Os requisitos de um sistema dividem-se, de forma clássica, em
requisitos funcionais, que descrevem os serviços que o
sistema deve oferecer, e requisitos não funcionais, que
estabelecem restrições de qualidade como segurança,
rastreabilidade, usabilidade e disponibilidade
\cite{Alves2022}. No contexto deste trabalho, os requisitos
funcionais contemplam a abertura padronizada de chamados, a
geração automática de protocolo, o controle de status e a
geração de indicadores gerenciais; os requisitos não
funcionais incluem a operação dentro do ecossistema Microsoft
365 já licenciado pela organização e o acesso diferenciado por
perfil de usuário \cite{Lopes2024}.

Entre as técnicas de elicitação de requisitos, o uso de
questionários e formulários estruturados destaca-se por
permitir a coleta padronizada de informações junto a
diferentes perfis de \textit{stakeholders}, com menor custo
operacional em comparação a técnicas presenciais
\cite{Gil2002}. A segmentação das perguntas conforme o papel
do respondente no processo --- nível operacional, de
supervisão ou de gestão --- favorece a triangulação das
informações coletadas, permitindo validar o mesmo fluxo sob
diferentes perspectivas organizacionais.

\subsection{Plataformas Low-Code como Estratégia Tecnológica}

As plataformas \textit{low-code} constituem uma estratégia
tecnológica para implementação de soluções corporativas,
especialmente em organizações que não dispõem de equipes de
desenvolvimento dedicadas. O paradigma \textit{low-code}
minimiza a necessidade de codificação manual por meio de
interfaces visuais e componentes pré-configurados, permitindo
que profissionais com menor especialização técnica construam
e mantenham aplicações funcionais \cite{Alves2022}.
\cite{Santos2023} destacam que essa característica tem
viabilizado o fenômeno dos \textit{citizen developers} ---
colaboradores das áreas de negócio que assumem papel ativo na
construção de soluções digitais sem depender exclusivamente
de equipes de TI --- e que a combinação entre BPM e
\textit{low-code} é especialmente eficaz na construção de
repositórios e painéis de acompanhamento, pois permite que as
regras de negócio sejam configuradas pelas próprias áreas
proprietárias dos processos. \cite{Costa2022} demonstraram,
em estudo de caso aplicado, que a adoção de ferramentas
\textit{low-code} resultou em rastreabilidade completa e
redução de 45\% no tempo de processamento de solicitações
internas, com desenvolvimento concluído em semanas.

Este trabalho não pretende apresentar o \textit{low-code}
como substituto da programação tradicional, pois não há, na
literatura consultada, métricas suficientes para sustentar
tal comparação. O foco adotado é demonstrar como a plataforma
possibilitou implementar uma solução integrada ao ambiente
corporativo já existente, reduzindo barreiras técnicas e
permitindo a criação de uma aplicação alinhada às
necessidades do processo redesenhado. A arquitetura adotada é
composta por quatro componentes do ecossistema Microsoft 365.
O primeiro é o \textbf{Power Apps}, interface de abertura e
acompanhamento do chamado. O segundo é a \textbf{Microsoft
List}, evolução das listas do SharePoint Online, que funciona
como repositório centralizado das informações, herdando
recursos de controle de versão, permissões por perfil e
integração nativa com o Power Automate e o Power BI
\cite{Mendes2024}. O terceiro é o \textbf{Power Automate},
responsável pela execução das regras de negócio e das
notificações automáticas. O quarto é o \textbf{Power BI},
responsável pela análise gerencial e geração de indicadores.
Essa arquitetura reflete o modelo de integração nativa entre
componentes do Microsoft 365, sem necessidade de conectores
externos ou infraestrutura adicional \cite{Costa2022}.

\subsection{Automação de Processos}

A automação de processos é o mecanismo pelo qual regras de
negócio previamente definidas passam a ser executadas de
forma padronizada, reduzindo a dependência de atividades
manuais para comunicação, direcionamento e acompanhamento de
demandas \cite{ReisFaria2024}. Estudos sobre automação
robótica de processos (RPA) em contextos organizacionais
distintos do estudado neste trabalho indicam ganhos
recorrentes de padronização e redução de retrabalho quando a
automação é aplicada sobre um processo previamente revisado,
e não sobre um fluxo manual simplesmente transposto para o
ambiente digital \cite{ReisFaria2024}.

No processo analisado, a automação substitui uma sequência de
etapas dependentes de ação humana --- identificação da
necessidade, envio de comunicação e acompanhamento manual do
retorno --- por uma sequência orientada a eventos, na qual a
alteração de um registro dispara automaticamente a regra
correspondente, a notificação ao responsável e a atualização
do histórico. Essa transformação está alinhada ao que
\cite{Lopes2024} descrevem como resultado da digitalização de
processos com ferramentas do ecossistema Microsoft: repositório
único de dados, com rastreabilidade completa e eliminação da
comunicação informal entre setores. A automação, portanto, não
é tratada neste trabalho como mecanismo isolado de registro,
mas como consequência do redesenho do processo discutido nas
seções anteriores.

\subsection{SLA na Gestão de Chamados Internos}

O Acordo de Nível de Serviço (\textit{Service Level
Agreement} --- SLA) é o instrumento que formaliza prazos,
critérios de qualidade e responsabilidades no atendimento de
demandas internas. \cite{Malheiros2024} concluem que o SLA é
fundamental para estabelecer expectativas claras entre
provedores e usuários de serviços, promovendo qualidade,
transparência e resolução eficaz de disputas, e destacam que
a automação do controle de prazos --- com geração de alertas
quando um chamado se aproxima do vencimento --- é condição
para que o SLA deixe de ser apenas um documento formal e passe
a ser um instrumento ativo de gestão. No contexto deste
trabalho, o SLA é incorporado como requisito não funcional e
como regra de negócio associada à prioridade de cada chamado,
integrando-se ao ciclo de monitoramento previsto pelo BPM.

\subsection{Comunicação Interna e Gestão de Informações entre
Áreas}

A comunicação interna é reconhecida na literatura como fator
crítico para a eficiência operacional e para o desempenho
organizacional. \cite{Souza2023} concluem, em revisão de
literatura, que grande parte das organizações tem dificuldade
em manter uma comunicação interna eficiente, optando, na
maioria das vezes, por canais que não garantem registro,
completude ou rastreabilidade das informações trocadas entre
setores. A substituição de canais informais por sistemas
estruturados, com registro automático e notificação
padronizada, reduz erros, elimina retrabalho e melhora o
desempenho das equipes envolvidas \cite{Lopes2024}.

No contexto específico da gestão de chamados internos, a
formalização do fluxo por meio de um canal estruturado
representa a passagem de um modelo reativo --- no qual o
problema é comunicado de forma \textit{ad hoc} --- para um
modelo proativo, no qual cada chamado gera automaticamente um
protocolo, uma notificação à área responsável e um registro
rastreável desde a abertura até o encerramento. Essa
transição é o elo final entre os seis eixos teóricos
apresentados nesta seção: o problema organizacional
compreendido pela Engenharia de Processos é redesenhado sob a
disciplina do BPM, especificado pela Engenharia de Requisitos,
implementado pela Engenharia de Software em ambiente
\textit{low-code} e operacionalizado pela automação, com a
comunicação estruturada entre áreas como resultado observável
da transformação.
